\documentclass[a4paper,12pt]{article}
\usepackage[brazil]{babel}
\usepackage[utf8]{inputenc}
\usepackage{geometry}
\usepackage{titlesec}
\usepackage{longtable}
\usepackage{fancyhdr}
\usepackage{graphicx}
\usepackage{float}
\usepackage{xcolor}
\usepackage{xurl}
\usepackage{hyperref}

\geometry{left=3cm, right=2cm, top=3cm, bottom=2cm}


\pagestyle{empty}
\pagenumbering{gobble} 

% Título e folha de rosto
\title{
    \vspace{4cm}
    \textbf{Fatec Zona Sul - Dom Paulo Evaristo Arns}\\
    \vspace{1cm}
    \textbf{Curso de Análise e Desenvolvimento de Sistemas}\\
    \vspace{3cm}
    \textbf{\large Uma Solução Prática para Avaliação de Performance em Interfaces de Programação de Aplicações
}\\
    \vspace{3cm}
    \textbf{Danillo Martins Leal}\\
    \textbf{Jean Marcelino de Lima}\\
    \vspace{3cm}
    São Paulo - SP\\
    2025
}
\author{}
\date{}

\newcommand{\folhaderosto}{
    \begin{titlepage}
        \centering
        \vspace*{2cm}

        \textbf{Danillo Martins Leal}\\
        \textbf{Jean Marcelino de Lima}\\[3cm]

        {\large \textbf{Uma Solução para Avaliação de Desempenho e Escalabilidade em Interfaces de Programação de Aplicações (APIs)}}\\[2cm]

        Trabalho de Conclusão de Curso apresentado à Faculdade de Tecnologia da Zona Sul – Dom Paulo Evaristo Arns, como requisito parcial para obtenção do título de Tecnólogo em Análise e Desenvolvimento de Sistemas.\\[1.5cm]

        Orientador: Prof. Paulo Rogério da Silva\\[3cm]

        São Paulo - SP\\
        2025

        \vfill
    \end{titlepage}
}

\begin{document}

\maketitle
\folhaderosto

\newpage
\tableofcontents
\newpage

\clearpage
\pagenumbering{arabic}
\setcounter{page}{4}
\pagestyle{plain}

% -------------------- Conteúdo ----------------------

\section{Introdução}
\subsection{Contextualização}
\subsection{Problemática}
\subsection{Objetivos}
\subsubsection{Objetivo Geral}
\subsubsection{Objetivos Específicos}
\subsection{Justificativa}
\subsection{Estrutura do Trabalho}

\section{Referencial Teórico}
\subsection{Testes de Performance e Carga}
\subsection{Conceitos de Escalabilidade, Latência e Throughput}
\subsection{Ferramentas Existentes para Teste de Carga}
\subsubsection{JMeter}
\subsubsection{k6}
\subsubsection{Gatling}
\subsubsection{Outras ferramentas}
\subsection{Tecnologias Utilizadas}
\subsubsection{Node.js / NestJS}
\subsubsection{React (se houver frontend)}
\subsubsection{Banco de Dados (PostgreSQL, MongoDB, etc.)}
\subsubsection{Outras ferramentas e bibliotecas}

\section{Metodologia}
\subsection{Tipo de Pesquisa}
\subsection{Procedimentos Metodológicos}
\subsection{Processos de Desenvolvimento}
\subsubsection{Levantamento de Requisitos}
\subsubsection{Desenvolvimento Incremental ou Metodologia Ágil}
\subsection{Ferramentas e Ambientes Utilizados}

\section{Modelagem do Sistema}
\subsection{Requisitos Funcionais e Não Funcionais}
\subsection{Diagrama de Casos de Uso}
\subsection{Diagrama de Classes}
\subsection{Diagrama de Sequência}
\subsection{Diagrama de Componentes}
\subsection{Diagrama de Implantação (opcional)}

\section{Desenvolvimento da Solução}
\subsection{Arquitetura da Aplicação}
\subsection{Backend – API e Executor de Testes}
\subsection{Frontend (se houver)}
\subsection{Banco de Dados}
\subsection{Funcionalidades Implementadas}
\subsubsection{Criação de Cenários de Teste}
\subsubsection{Execução dos Testes de Carga}
\subsubsection{Geração de Relatórios e Métricas}
\subsection{Desafios e Soluções Encontradas}

\section{Resultados e Discussões}
\subsection{Execução dos Testes com a Ferramenta Desenvolvida}
\subsection{Análise dos Resultados}
\subsection{Comparação com Ferramentas Existentes (opcional)}
\subsection{Limitações da Solução}

\section{Conclusão}
\subsection{Considerações Finais}
\subsection{Contribuições do Trabalho}
\subsection{Trabalhos Futuros}

\section{Referências}

\section{Apêndices e Anexos}
% Inserir se necessário

\end{document}
